\documentclass[12pt,t,handout]{beamer}
\usepackage{graphicx}
\setbeameroption{show notes}
\setbeamertemplate{note page}[plain]
\usepackage{listings}
\usepackage{datetime}
\usepackage{url}

% specifications for presenter mode
%\beamerdefaultoverlayspecification{<+->}
%\setbeamercovered{transparent}

\usepackage[english]{babel}
\usepackage[utf8x]{inputenc}

\usepackage{amstext}
%\usepackage{coloremoji}

\usepackage{graphicx}
\graphicspath{ {Figs/} }

% math shorthand
\usepackage{bm}
\usepackage{amsmath}
\usepackage{mathtools}
\newcommand{\D}{\mathcal{D}}
\newcommand{\E}{\mathbb{E}}
\newcommand{\F}{\mathcal{F}}
\newcommand{\X}{\mathcal{X}}
\newcommand{\lik}{\mathcal{L}}
\DeclarePairedDelimiterX{\infdivx}[2]{(}{)}{%
  #1\;\delimsize\|\;#2%
}
\newcommand{\infdiv}{D\infdivx}
\DeclarePairedDelimiter{\norm}{\lVert}{\rVert}
\DeclareMathOperator*{\argmin}{arg\,min}
\DeclareMathOperator*{\argmax}{arg\,max}

% Bibliography
\usepackage{natbib}
\bibpunct{(}{)}{,}{a}{}{;}
\usepackage{bibentry}
\nobibliography*

\def\notescolors{1}
% header.tex: boring LaTeX/Beamer details + macros

% get rid of junk
\usetheme{default}
\beamertemplatenavigationsymbolsempty
\hypersetup{pdfpagemode=UseNone} % don't show bookmarks on initial view


% font
\usepackage{fontspec}
\setsansfont
  [ ExternalLocation = fonts/ ,
    UprightFont = *-regular ,
    BoldFont = *-bold ,
    ItalicFont = *-italic ,
    BoldItalicFont = *-bolditalic ]{texgyreheros}
\setbeamerfont{note page}{family*=pplx,size=\footnotesize} % Palatino for notes
% "TeX Gyre Heros can be used as a replacement for Helvetica"
% I've placed them in fonts/; alternatively you can install them
% permanently on your system as follows:
%     Download http://www.gust.org.pl/projects/e-foundry/tex-gyre/heros/qhv2.004otf.zip
%     In Unix, unzip it into ~/.fonts
%     In Mac, unzip it, double-click the .otf files, and install using "FontBook"

% named colors
\definecolor{offwhite}{RGB}{255,250,240}
\definecolor{gray}{RGB}{155,155,155}

\ifx\notescolors\undefined % slides
  \definecolor{foreground}{RGB}{255,255,255}
  \definecolor{background}{RGB}{24,24,24}
  \definecolor{title}{RGB}{107,174,214}
  \definecolor{subtitle}{RGB}{102,255,204}
  \definecolor{hilit}{RGB}{102,255,204}
  \definecolor{vhilit}{RGB}{255,111,207}
  \definecolor{lolit}{RGB}{155,155,155}
\else % notes
  \definecolor{background}{RGB}{255,255,255}
  \definecolor{foreground}{RGB}{24,24,24}
  \definecolor{title}{RGB}{27,94,134}
  \definecolor{subtitle}{RGB}{22,175,124}
  \definecolor{hilit}{RGB}{122,0,128}
  \definecolor{vhilit}{RGB}{255,0,128}
  \definecolor{lolit}{RGB}{95,95,95}
\fi
\definecolor{nhilit}{RGB}{128,0,128}  % hilit color in notes
\definecolor{nvhilit}{RGB}{255,0,128} % vhilit for notes

\newcommand{\hilit}{\color{hilit}}
\newcommand{\vhilit}{\color{vhilit}}
\newcommand{\nhilit}{\color{nhilit}}
\newcommand{\nvhilit}{\color{nvhilit}}
\newcommand{\lolit}{\color{lolit}}

% use those colors
\setbeamercolor{titlelike}{fg=title}
\setbeamercolor{subtitle}{fg=subtitle}
\setbeamercolor{institute}{fg=lolit}
\setbeamercolor{normal text}{fg=foreground,bg=background}
\setbeamercolor{item}{fg=foreground} % color of bullets
\setbeamercolor{subitem}{fg=lolit}
\setbeamercolor{itemize/enumerate subbody}{fg=lolit}
\setbeamertemplate{itemize subitem}{{\textendash}}
\setbeamerfont{itemize/enumerate subbody}{size=\footnotesize}
\setbeamerfont{itemize/enumerate subitem}{size=\footnotesize}

% page number
\setbeamertemplate{footline}{%
    \raisebox{5pt}{\makebox[\paperwidth]{\hfill\makebox[20pt]{\lolit
          \scriptsize\insertframenumber}}}\hspace*{5pt}}

% add a bit of space at the top of the notes page
\addtobeamertemplate{note page}{\setlength{\parskip}{12pt}}

% default link color
\hypersetup{colorlinks, urlcolor={hilit}}

\ifx\notescolors\undefined % slides
  % set up listing environment
  \lstset{language=bash,
          basicstyle=\ttfamily\scriptsize,
          frame=single,
          commentstyle=,
          backgroundcolor=\color{darkgray},
          showspaces=false,
          showstringspaces=false
          }
\else % notes
  \lstset{language=bash,
          basicstyle=\ttfamily\scriptsize,
          frame=single,
          commentstyle=,
          backgroundcolor=\color{offwhite},
          showspaces=false,
          showstringspaces=false
          }
\fi

% a few macros
\newcommand{\bi}{\begin{itemize}}
\newcommand{\bbi}{\vspace{24pt} \begin{itemize} \itemsep8pt}
\newcommand{\ei}{\end{itemize}}
\newcommand{\ig}{\includegraphics}
\newcommand{\subt}[1]{{\footnotesize \color{subtitle} {#1}}}
\newcommand{\ttsm}{\tt \small}
\newcommand{\ttfn}{\tt \footnotesize}
\newcommand{\figh}[2]{\centerline{\includegraphics[height=#2\textheight]{#1}}}
\newcommand{\figw}[2]{\centerline{\includegraphics[width=#2\textwidth]{#1}}}


%%%%%%%%%%%%%%%%%%%%%%%%%%%%%%%%%%%%%%%%%%%%%%%%%%%%%%%%%%%%%%%%%%%%%%
% end of header
%%%%%%%%%%%%%%%%%%%%%%%%%%%%%%%%%%%%%%%%%%%%%%%%%%%%%%%%%%%%%%%%%%%%%%

% title info
\title{\large Efficient Nonparametric Inference for Stochastic Interventions
  Under Censoring}
\subtitle{\scriptsize for the UC Berkeley Biostatistics Seminar Series,\\ given
                      02 April 2018 \\[-10pt]
         }
\author{\href{https://nimahejazi.org}{Nima Hejazi}
       \\[-10pt]
       }
\institute{Group in Biostatistics \\
           University of California, Berkeley \\
           \href{https://www.stat.berkeley.edu/~nhejazi}
             {\tt \scriptsize \color{foreground}
               stat.berkeley.edu/\textasciitilde{}nhejazi
             }
           \\[4pt]
           \includegraphics[height=20mm]{Figs/seal-berkeley.png}
           \\[-12pt]
          }
\date{
  \href{https://nimahejazi.org}
      {\tt \scriptsize \color{foreground} nimahejazi.org}
  \\[-4pt]
  \href{https://twitter.com/nshejazi}
      {\tt \scriptsize \color{foreground} twitter/@nshejazi}
  \\[-4pt]
  \href{https://github.com/nhejazi}
      {\tt \scriptsize \color{foreground} github/nhejazi}
}

%%%%%%%%%%%%%%%%%%%%%%%%%%%%%%%%%%%%%%%%%%%%%%%%%%%%%%%%%%%%%%%%%%%%%%%%%%%%%%%%

\begin{document}

% title slide
{
\setbeamertemplate{footline}{} % no page number here
\frame{
  \titlepage

  \vspace{-1em}

  \centerline{\href{https://goo.gl/xabp3Q}{\tt \scriptsize
                                           \underline{slides}: goo.gl/xabp3Q}}
  \vspace{-1.5em}
  \vfill \hfill \includegraphics[height=6mm]{Figs/cc-zero.png} \vspace*{-0.5cm}

  \note{This slide deck is for a seminar-length talk (about 50 minutes) on a new
    approach to causal inference and nonparametric variable importance in the
    context of parameters defined as treatment shifts. Here, we introduce an
    additive treatment shift parameter, extensions for censored data (including
    a multiple double robustness property), new statistical software for
    applying our approach, and applications to a vaccine efficacy trial
    examining HIV. This talk was most recently given at a meeting of the
    Biostatistics Seminar Series at the University of California, Berkeley.

    Source: {\tt https://github.com/nhejazi/talk\_txshift} \\
    Slides: {\tt https://goo.gl/JDhSEg} \\
    With notes: {\tt https://goo.gl/xabp3Q}
}
}
}

%%%%%%%%%%%%%%%%%%%%%%%%%%%%%%%%%%%%%%%%%%%%%%%%%%%%%%%%%%%%%%%%%%%%%%%%%%%%%%%%

\begin{frame}[c]{Preview: Summary}
\only<1>{\addtocounter{framenumber}{-1}}

\begin{center}
\begin{itemize}
  \itemsep10pt
  \item ...
  \item ...
  \item ...
  \item ...
\end{itemize}
\end{center}

\note{We'll go over this summary again at the end of the talk. Hopefully, it
  will all make more sense then.
}

\end{frame}

%%%%%%%%%%%%%%%%%%%%%%%%%%%%%%%%%%%%%%%%%%%%%%%%%%%%%%%%%%%%%%%%%%%%%%%%%%%%%%%%

\begin{frame}[c]{Motivation: Let's meet the data}

\begin{center}
\begin{itemize}
  \itemsep10pt
  \item Observational study of the impact of disease state on DNA methylation.
  \item Phenotype-level quantities: $216$ subjects, binary disease status (FASD)
    of each subject, background info on subjects (e.g., sex, age).
  \item Genomic-level quantities: $\sim 850,000$ CpG sites interrogated using
    the \textit{Infinium MethylationEPIC BeadChip} by Illumina.
  \item \textbf{Questions}: How do disease status and differential methylation
    relate? Is a coherent biomarker-type signature detectable?
\end{itemize}
\end{center}

\note{
\begin{itemize}
  \itemsep10pt
  \item FASD is an abbreviation for Fetal Alcohol Spectrum Disorders.
  \item We're mostly interested in the interplay between disease and DNA
    methylation.
  \item In particular, we'd like to construct some kind of importance score for
    CpG sites impacted by the exposure/disease of interest.
  \item Re: dimensionality, c.f., RNA-seq analyses are $\sim 30,000$ in
    dimension at the gene level.
\end{itemize}
}

\end{frame}

%%%%%%%%%%%%%%%%%%%%%%%%%%%%%%%%%%%%%%%%%%%%%%%%%%%%%%%%%%%%%%%%%%%%%%%%%%%%%%%%

\begin{frame}[c]{DNA Methylation}

%\begin{figure}[H]
  %\centering
  %\includegraphics[width=\textwidth]{dna_methylation}
  %\caption{
    %\url{https://www.illumina.com/techniques/sequencing/methylation-sequencing.html}
    %(source)}
%\end{figure}

\note{
\begin{itemize}
  \itemsep10pt
  \item The biology of these structures is quite complicated: many different
    structures --- CpG islands, shores, etc.
  \item Both hypermethylation and hypomethylation are linked to disease states.
  \item It's especially important to examine these processes, keeping in mind
    that the process may be complex and non-monotonic.
\end{itemize}
}

\end{frame}

%%%%%%%%%%%%%%%%%%%%%%%%%%%%%%%%%%%%%%%%%%%%%%%%%%%%%%%%%%%%%%%%%%%%%%%%%%%%%%%%

\begin{frame}[c]{Data analysis? Linear Models!}
\begin{center}
\begin{itemize}
  \itemsep10pt
  \item Standard operating procedure: For each CpG site ($g = 1, \dots, G$), fit
    a linear model:
    \[
      \mathbb{E}[y_g] = X \beta_g
    \]
  \item Test the coefficent of interest using a standard t-test:
    \[
      t_{g} = \frac{\hat{\beta}_{g} - \beta_{g, H_0}}{s_g}
    \]
  \item Such models are a matter of convenience: does $\hat{\beta}_{g}$ answer
    our scientific questions? Perhaps not.
  \item Is consideration being given to whether the data could have been
    generated by a linear model? Perhaps not.
\end{itemize}
\end{center}

\note{
\begin{itemize}
  \itemsep10pt
  \item CpG sites are thought to function in networks. Treating them as acting
    independently is \textit{not} faithful to the underlying biology.
  \item The linear model is a great starting point for analyses whne the data is
    generated using complex technology --- no need to make the analysis more
    complicated.
  \item That being said, the data is difficult and expensive to collect, so why
    restrict the scope of the questions we'd like to ask.
\end{itemize}
}
\end{frame}

%%%%%%%%%%%%%%%%%%%%%%%%%%%%%%%%%%%%%%%%%%%%%%%%%%%%%%%%%%%%%%%%%%%%%%%%%%%%%%%%%%%%

\begin{frame}[fragile,c]{Motivation: Science Before Statistics}

\begin{center}
\begin{minipage}[c]{9.4cm}
\begin{semiverbatim}
\lstset{basicstyle=\normalsize}
\begin{lstlisting}[linewidth=9.4cm]
What is the effect of disease status
on DNA methylation at a specific CpG
site, controlling for the observed
methylation status of the neighbors
of the given CpG site?

\end{lstlisting}
\end{semiverbatim}
\end{minipage}
\end{center}

\note{
\begin{itemize}
  \itemsep10pt
  \item Again, CpG sites are thought to function in networks. Treating them as
    acting independently is \textit{not} faithful to the underlying biology.
  \item This means that we should take into account the methylation status of
    neighboring CpG sites when assessing differential methylation at a single
    site.
  \item This is a coherent scientific question that we can set out to answer
    statistically. It's motivated by the established science and possible to do
    with modern statistical methodology.
\end{itemize}
}

\end{frame}



%%%%%%%%%%%%%%%%%%%%%%%%%%%%%%%%%%%%%%%%%%%%%%%%%%%%%%%%%%%%%%%%%%%%%%%%%%%%%%%%%%%%

\begin{frame}[c]{Data analysis? A Data-Adaptive Approach}

\begin{center}
\begin{enumerate}
  \itemsep8pt
  \item Isolate a subset of CpG sites for which there is cursory evidence of
    differential methylation.
  \item Assign CpG sites into neighborhoods (e.g., bp distance). If there are
    many neighbors, apply clustering (e.g., \textit{PAM}) to select a subset.
  \item Estimate \textit{variable importance measure} (VIM) at each screened CpG
    site, with disease as intervention ($A$) and controlling for neighboring CpG
    sites ($W$).
  \item Apply a variant of the Benjamini \& Hochberg method for FDR control,
    accounting for initial screening.
\end{enumerate}
\end{center}

\note{
\begin{itemize}
  \itemsep10pt
  \item Pre-screening is a critical step since we cannot perform computationally
    intensive estimation on all the sites. This is flexible --- just use your
    favorite method (as long as allows a ranking to be made).
  \item The variable importance step merely comes down to the creation of a
    score. We use TMLE to statistically estimate parameters from causal models.
    The procedure is general enough to accomodate any inference technique.
\end{itemize}
}

\end{frame}

%%%%%%%%%%%%%%%%%%%%%%%%%%%%%%%%%%%%%%%%%%%%%%%%%%%%%%%%%%%%%%%%%%%%%%%%%%%%%%%%

\begin{frame}[c]{Pre-Screening --- Pick Your Favorite Method}

\begin{center}
\begin{itemize}
  \itemsep10pt
  \item The estimation procedure is computationally intensive --- apply it only
    to sites that appear promising.
  \item Consider estimating univariate (linear) regressions of intervention on
    CpG methylation status. Fast, easy.
  \item Select CpG sites with a marginal p-value below, say, $0.01$. Apply
    data-adaptive procedure to this subset.
  \item The modeling assumptions do not matter since the we won't be pursuing
    inference under such a model.
  \item Software implementation is extensible. Users are encouraged to add their
    own. (It's easy!)
\end{itemize}
\end{center}

\note{
\begin{itemize}
  \itemsep10pt
  \item We'll be adding to the available routines for pre-screening too! For
    now, we have \textsc{limma}, and more are on the way.
\end{itemize}
}

\end{frame}

%%%%%%%%%%%%%%%%%%%%%%%%%%%%%%%%%%%%%%%%%%%%%%%%%%%%%%%%%%%%%%%%%%%%%%%%%%%%%%%%

\begin{frame}[c]{Too Many Neighbors? Clustering}

\begin{center}
\begin{itemize}
  \itemsep10pt
  \item There are many options: $k$-means, $k$-medoids, etc., as well as many
    algorithmic solutions.
  \item For convenience, we use \underline{P}artitioning \underline{A}round
    \underline{M}edoids (PAM), a well-established algorithm.
  \item With limited sample sizes, the number of neighboring sites that may be
    controlled for is limited.
  \item To faithfully answer the question of interest, choose the neighboring
    sites that are the most representative.
  \item This is an \textit{optional} step --- it need only be applied when there
    is a large number of CpG sites in the neighborhood of the target CpG site.
\end{itemize}
\end{center}

\note{
\begin{itemize}
  \itemsep10pt
  \item The number of sites that we can control for is roughly a function of
    sample size. This impacts the definition of the parameter that we estimate,
    and allows enough flexibility to obtain either very local or more regional
    estimates.
\end{itemize}
}

\end{frame}

%%%%%%%%%%%%%%%%%%%%%%%%%%%%%%%%%%%%%%%%%%%%%%%%%%%%%%%%%%%%%%%%%%%%%%%%%%%%%%%%

\begin{frame}[c]{Nonparametric Variable Importance}

\begin{center}
\begin{itemize}
  \itemsep8pt
  \item Let's consider a simple target parameter: the average treatment effect
    (ATE):
    \[
      \Psi_g(P_0) = \mathbb{E}_{W,0}[\mathbb{E}_0[Y_g \mid A = 1, W_{-g}] -
      \mathbb{E}_0[Y_g \mid A = 0, W_{-g}]]
    \]
  \item Under certain (untestable) assumptions, interpretable as difference in
    methylation at site $g$ with intervention and, possibly contrary to
    fact, the same under no intervention, controlling for neighboring sites.
  \item Provides a \textit{nonparametric} (model-free) measure for those CpG
    sites impacted by a discrete intervention.
  \item Let the choice of parameter be determined by our scientific question of
    interest.
\end{itemize}
\end{center}

\note{
By allowing scientific questions to inform the parameters that we choose
to estimate, we can do a better job of actually answering the questions of
interest to our collaborators. Further, we abandon the need to specify the
functional relationship between our outcome and covariates; moreover, we
can now make use of advances in machine learning.
}

\end{frame}

%%%%%%%%%%%%%%%%%%%%%%%%%%%%%%%%%%%%%%%%%%%%%%%%%%%%%%%%%%%%%%%%%%%%%%%%%%%%

\begin{frame}[c]{Target Minimum Loss-Based Estimation}

\begin{center}
\begin{itemize}
  \itemsep8pt
  \item We use \textit{targeted minimum loss-based estimation} (TMLE), a method
    for inference in semiparametric infinite-dimensional statistical models.
  \item No need to specify a functional form or assume that we know the true
    data-generating distribution.
  \item \textbf{Asymptotic linearity:}
    \[
      \Psi_g(P_n^*) - \Psi_g(P_0) = \frac{1}{n} \sum_{i = 1}^{n} IC(O_i) +
      o_P\left(\frac{1}{\sqrt{n}}\right)
    \]
  \item \textbf{Limiting distribution:}
    \[
      \sqrt{n}(\Psi_n - \Psi) \to N(0, Var(D(P_0)))
    \]
  \item \textbf{Statistical inference:}
    \[
      \Psi_n \pm z_{\alpha} \cdot \frac{\sigma_n}{\sqrt{n}}
    \]
\end{itemize}
\end{center}

\note{
Under the additional condition that the remainder term $R(\hat{P}^*, P_0)$
decays as $o_P \left( \frac{1}{\sqrt{n}} \right),$ we have that
$\Psi_n - \Psi_0 = (P_n - P_0) \cdot D(P_0) + o_P
\left( \frac{1}{\sqrt{n}} \right),$ which, by a central limit theorem,
establishes a Gaussian limiting distribution for the estimator, with variance
$V(D(P_0))$, the variance of the efficient influence curve (canonical gradient)
when $\Psi$ admits an asymptotically linear representation.

The above implies that $\Psi_n$ is a $\sqrt{n}$-consistent estimator of $\Psi$,
that it is asymptotically normal (as given above), and that it is locally
efficient. This allows us to build Wald-type confidence intervals, where
$\sigma_n^2$ is an estimator of $V(D(P_0))$. The estimator $\sigma_n^2$
may be obtained using the bootstrap or computed directly via $ \sigma_n^2 =
\frac{1}{n} \sum_{i = 1}^{n} D^2(\bar{Q}_n^*, g_n)(O_i)$
}

\end{frame}

%%%%%%%%%%%%%%%%%%%%%%%%%%%%%%%%%%%%%%%%%%%%%%%%%%%%%%%%%%%%%%%%%%%%%%%%%%%%%%%%

\begin{frame}[c]{Corrections for Multiple Testing}

\begin{center}
\begin{itemize}
  \itemsep10pt
  \item Multiple testing corrections are critical. Without these, we
    systematically obtain misleading results.
  \item The Benjamini \& Hochberg procedure for controlling the False Discovery
    Rate (FDR) is a well-established technique for addressing the multiple
    testing issue.
  \item We use a modified BH-FDR procedure to account for the pre-screening step
    of the proposed algorithm.
  \item This modified BH-FDR procedure for multi-stage analyses (FDR-MSA) works
    by adding a p-value of $1.0$ for each site that did not pass pre-screening
    then performs BH-FDR as normal.
\end{itemize}
\end{center}

\note{
\begin{itemize}
  \itemsep10pt
  \item Note that $\text{FDR} = \E\left[\frac{V}{R}\right] =
    \E\left[\frac{V}{R} \mid R > 0 \right] P(R > 0)$.
  \item BH-FDR procedure: Find $\hat{k} = max\{k: p_{(k)} \leq \frac{k}{M} \cdot
    \alpha\}$
  \item FDR-MSA will only incur a loss of power if the initial screening
    excludes variables that would have been rejected by the BH procedure when
    applied to the subset on which estimation was performed.
  \item BH-FDR control is a rank-based procedure, so we must assume that the
    pre-screening does not disrupt the ranking with respect to the estimation
    subset, which is provably true for screening procedures of a given type.
  \item MSA controls type I error with any procedure that is a function of only
    the type I error itself --- e.g., FWER. This does not hold for the FDR in
    complete generality.
\end{itemize}
}

\end{frame}

%%%%%%%%%%%%%%%%%%%%%%%%%%%%%%%%%%%%%%%%%%%%%%%%%%%%%%%%%%%%%%%%%%%%%%%%%%%%%%%%

\begin{frame}[c]{Software package: R/txshift}

%\begin{figure}[H]
  %\centering
  %\includegraphics[width=\textwidth]{methyvim_bioc}
  %\caption{
    %\url{https://github.com/nhejazi/txshift}
  %}
%\end{figure}

\begin{center}
\begin{itemize}
  \itemsep4pt
  \item Variable importance for continuous interventions.
  \item Take it for a test drive!
\end{itemize}
\end{center}

\note{
\begin{itemize}
  \itemsep10pt
  \item Contribute on GitHub.
  \item Reach out to us with questions and feature requests.
\end{itemize}
}

\end{frame}

%%%%%%%%%%%%%%%%%%%%%%%%%%%%%%%%%%%%%%%%%%%%%%%%%%%%%%%%%%%%%%%%%%%%%%%%%%%%%%%%

\begin{frame}[c]{Review: Summary}

\begin{center}
\begin{itemize}
  \itemsep10pt
  \item DNA methylation data is \textit{extremely} high-dimensional --- we can
    collect data on 850K genomic sites with modern arrays!
  \item Normalization and QC are critical components of properly analyzing
    modern DNA methylation data. There are many choices of technique.
  \item A relative scarcity of techniques for estimation and inference exists
    --- analyses are often limited to the general linear model.
  \item Statistical causal inference provides an avenue for answering richer
    scientific questions, especially when combined with modern advances in
    machine learning.
\end{itemize}
\end{center}

\note{It's always good to include a summary.}

\end{frame}

%%%%%%%%%%%%%%%%%%%%%%%%%%%%%%%%%%%%%%%%%%%%%%%%%%%%%%%%%%%%%%%%%%%%%%%%%%%%%%%%

% don't want dimming with references
\setbeamercovered{}
\beamerdefaultoverlayspecification{}

\begin{frame}[c,allowframebreaks]{References}

\bibliographystyle{apalike}
\nocite{*}
\bibliography{references}

%\note{Here's some work we've talked about. Go check these out if interested.}

\end{frame}

%%%%%%%%%%%%%%%%%%%%%%%%%%%%%%%%%%%%%%%%%%%%%%%%%%%%%%%%%%%%%%%%%%%%%%%%%%%%%%%%

\begin{frame}{Acknowledgments}

\vspace{20pt}

\underline{Collaborators:}

\begin{tabular}{@{}l@{\hspace{1.5cm}}l@{}}
  David Benkeser & \footnotesize \lolit Emory University \\
  Mark van der Laan & \footnotesize \lolit University of California, Berkeley \\
  Peter Gilbert & \footnotesize \lolit Fred Hutchinson Cancer Research Center \\
  Holly Janes & \footnotesize \lolit Fred Hutchinson Cancer Research Center \\
\end{tabular}

\vspace{10mm}

\underline{Funding source}:\\
NIH Biomedical Big Data Training Program: T32-LM012417-02

\end{frame}

%%%%%%%%%%%%%%%%%%%%%%%%%%%%%%%%%%%%%%%%%%%%%%%%%%%%%%%%%%%%%%%%%%%%%%%%%%%%%%%%

\begin{frame}[c]{Thank you.}

\Large
Slides: \href{https://goo.gl/JDhSEg}{goo.gl/JDhSEg} \quad
\includegraphics[height=5mm]{Figs/cc-zero.png}

\vspace{3mm}
Notes: \href{https://goo.gl/xabp3Q}{goo.gl/xabp3Q}

\vspace{3mm}
Source (repo): \href{https://goo.gl/m5As73}{goo.gl/m5As73}

\vspace{3mm}
\href{https://www.stat.berkeley.edu/~nhejazi}{\tt
  stat.berkeley.edu/\textasciitilde{}nhejazi}

\vspace{3mm}
\href{https://nimahejazi.org}{\tt nimahejazi.org}

\vspace{3mm}
\href{https://twitter.com/nshejazi}{\tt twitter/@nshejazi}

\vspace{3mm}
\href{https://github.com/nhejazi}{\tt github/nhejazi}

%\note{Here's where you can find me, as well as the slides for this talk.}

\end{frame}

%%%%%%%%%%%%%%%%%%%%%%%%%%%%%%%%%%%%%%%%%%%%%%%%%%%%%%%%%%%%%%%%%%%%%%%%%%%%%%%%

\end{document}

